% =============================================================================
% Module 21: Solutions to Exercises
% =============================================================================
% This file is conditionally included from main.tex via \ifsolutions
% =============================================================================

\section{Solutions to Exercises}
\label{sec:pe_solutions}

\noindent
Full worked solutions are also available in the Mathematica script
\solnlink{21_Perturbative_Estimators/problems\_21.wl}{problems\_21.wl},
which verifies each answer symbolically and, where a number is quoted,
numerically. Throughout, $\thetaE=1$, $r_j=\mBH a_j-\epsilon b_j-\delta\ysrc$
with $a_j=1/\theta_j$ and $b_j=(\gp-1)\alpha_{\rm PL}(\theta_j)$, and
$\epsilon=\mBH c$ with $c=\operatorname{Cov}_w(a,b)/\operatorname{Var}_w(b)$.

% ----- Exercise: two images -----
\subsection*{Exercise~\ref{ex:pe_two_images}: Two Images}

With two images the system
\[
  \mBH a_1-\epsilon b_1-\delta\ysrc=0,\qquad
  \mBH a_2-\epsilon b_2-\delta\ysrc=0
\]
has two equations for the two unknowns $(\epsilon,\delta\ysrc)$.
Subtracting eliminates $\delta\ysrc$:
\[
  \mBH(a_1-a_2)=\epsilon\,(b_1-b_2)
  \quad\Longrightarrow\quad
  \epsilon=\mBH\,\frac{a_1-a_2}{b_1-b_2},
\]
so $c=(a_1-a_2)/(b_1-b_2)$, and back-substitution gives
$\delta\ysrc=\mBH a_1-\epsilon b_1$. The general regression formula
reduces to the same difference quotient for $n=2$.

The residual vanishes because a square system that is
non-singular ($b_1\neq b_2$, i.e.\ the two images are not at the same
signed position) has an exact solution, so the least-squares minimum is
zero. With three images there are three equations for the same two
unknowns; the point-mass shape $a_j\propto1/\theta_j$ is generically
not a linear combination of $b_j$ and a constant, and no choice
of $(\epsilon,\delta\ysrc)$ annihilates all three residuals. For the
canonical triple at $\rho=0.1$, $\gp=1.5$ the minimum residual is
strictly positive; it vanishes only in the degenerate limit $\rho=1$.

The leftover residual is the only part of the central mass's signal
that the model cannot absorb: it is the evidence that a central mass is
present. A plane with
two images can be biased by a central mass but can never detect one; a
plane with a radial image can, if the slope is held fixed. With the
slope free, a single three-image plane gives three equations for three
unknowns and is fitted exactly (\S\ref{sec:pe_free_slope}; for the
running example the $3\times3$ system has determinant $0.561$), so
detection then needs a fourth datum or a second source plane that shares
the slope.
\emph{[Attested: \solnlink{21_Perturbative_Estimators/problems\_21.wl}{problems\_21.wl}, checks 1a--1f.]}

% ----- Exercise: intercept -----
\subsection*{Exercise~\ref{ex:pe_intercept}: Why the Intercept Matters}

The response coefficient is built from centred weighted moments,
\[
  c=\frac{\sum_jw_j(a_j-\bar a)(b_j-\bar b)}{\sum_jw_j(b_j-\bar b)^2}.
\]
Under $a_j\to a_j+\kappa$ the weighted mean shifts by the same constant,
$\bar a\to\bar a+\kappa$, so every $a_j-\bar a$ is unchanged and so is
$c$. The intercept absorbs the shift: from
$\delta\ysrc=\mBH(\bar a-c\bar b)$, $\delta\ysrc\to\delta\ysrc+\mBH\kappa$,
while $\epsilon$ is unchanged.

The physical content is that a constant added to every deflection on a
plane is exactly a uniform shift of that plane's source, and the source
position is not observable. The invariance says that the estimator uses
only observable information.

If the intercept is fixed, setting $\delta\ysrc=0$ and regressing
through the origin gives $c_{\rm fixed}=\sum_ja_jb_j/\sum_jb_j^2$. For
the canonical triple at $\rho=0.1$, $\gp=1.5$ this differs from
Eq.~\eqref{eq:pe_cclosed}. Fixing the source would let a pure source
shift masquerade as information about the Einstein radius.
\emph{[Attested: \solnlink{21_Perturbative_Estimators/problems\_21.wl}{problems\_21.wl}, checks 2a--2c.]}

% ----- Exercise: radial leverage -----
\subsection*{Exercise~\ref{ex:pe_radial_leverage}: The Radial Image's Leverage}

(a) The radial image enters $a$ as $1/\rho\to\infty$ (the
point-mass deflection diverges at small radius) but enters $b$ as
$(\gp-1)\rho^{2-\gp}\to0$ (the power law's response to $\thetaE$ is
weakest there). The uncentred cross moment is
$\sum_ja_jb_j=(\gp-1)[2+\rho^{1-\gp}]$: the pair contributes $2$, the
radial image $\rho^{1-\gp}\to\infty$. After centring,
$\operatorname{Cov}(a,b)\to\tfrac23(\gp-1)\rho^{1-\gp}$, still carried by
the radial image, while $\operatorname{Var}(b)\to2(\gp-1)^2$ is set by
the pair. Hence
\[
  c\simeq\frac{\tfrac23(\gp-1)\rho^{1-\gp}}{2(\gp-1)^2}
  =\frac{\rho^{1-\gp}}{3(\gp-1)}\xrightarrow[\rho\to0]{}\infty .
\]

(b) Nothing in (a) required $\rho\ll1$ until the last step.
The pair cancels in both means, so exactly $\bar a=1/(3\rho)$ and
$\bar b=(\gp-1)\rho^{2-\gp}/3$, and
\[
  \operatorname{Cov}(a,b)=2(\gp-1)\Bigl[1+\tfrac13\rho^{1-\gp}\Bigr],
  \qquad
  \operatorname{Var}(b)=2(\gp-1)^2\Bigl[1+\tfrac13\rho^{4-2\gp}\Bigr],
\]
whose ratio is Eq.~\eqref{eq:pe_cclosed},
$c=(3+\rho^{1-\gp})/[(\gp-1)(3+\rho^{4-2\gp})]$.

(c) Dividing by the asymptote,
$c/c_{\rm asym}=(1+3u)/(1+v/3)$ with $u=\rho^{\gp-1}$ and
$v=\rho^{4-2\gp}$. The term $3u$ comes from the ``$1$'' in the
covariance bracket, i.e.\ from the tangential pair; the term $v/3$ comes
from the variance bracket, i.e.\ from the radial image's $b_C^2$. For
small $\rho$ the term with the smaller exponent dominates, and
$\gp-1<4-2\gp\iff\gp<5/3$: below $\gp=5/3$ the pair's $+3u$ leads and
$c>c_{\rm asym}$; above it the radial image's $v/3$ leads and $c$ dips
below the asymptote at small $\rho$.

(d) At $\rho=1$: $c=(3+1)/[(\gp-1)(3+1)]=1/(\gp-1)$, the
two-image floor. The radial image has merged with a tangential one.

(e) At $\rho=0.03$, $\gp=1.4$:
\[
  c_{\rm asym}=3.388,\qquad
  c_{\rm one\ term}=5.871,\qquad
  c_{\rm exact}=5.859,
\]
and a direct numerical regression on $\theta=(1,-1,0.03)$ also gives
$5.859$. The asymptote is $42\%$ low (the true value is $1.73$ times
it), an error that would propagate unchanged into any bias built on
$c$; the one-term form is high by only $0.2\%$. The exact
Eq.~\eqref{eq:pe_cclosed} is the one to carry; the one-term form is
quoted only when a formula with visible structure is wanted. For a real lens with an asymmetric
pair, neither replaces the general regression on the actual images
(\S\ref{sec:pe_beyond_canonical}).
\emph{[Attested: \solnlink{21_Perturbative_Estimators/problems\_21.wl}{problems\_21.wl}, checks 3a--3k.]}

% ----- Exercise: sign change -----
\subsection*{Exercise~\ref{ex:pe_sign_change}: Sign Change}

From $\Delta\Brat/\Brat=(\gp-1)\mBH_1[c_1-\Brat^{k}c_2]$ the bias
vanishes, and changes sign, where the bracket does:
\[
  \frac{c_2}{c_1}=\Brat^{-k},\qquad k=\frac{3-\gp}{\gp-1}.
\]
For $1<\gp<3$, $k>0$, and for $0<\Brat<1$ (write $\Brat=e^{-L}$, $L>0$)
the threshold $\Brat^{-k}=e^{kL}>1$. The sign flips when plane~2's
response exceeds plane~1's by more than the geometric suppression
factor. Since $\partial(\Delta\Brat/\Brat)/\partial c_2=-(\gp-1)\mBH_1\Brat^k<0$,
a larger $c_2$ drives the bias negative; in a source-plane fit a larger
$c_2$ means a deeper plane-2 radial image, while in an image-plane fit
a very deep radial image is down-weighted and $c_2$ returns towards the
floor (\S\ref{sec:pe_beyond_canonical}). The sign of the cosmological
bias is controlled by where the radial images sit and by how the fit
weights them.

The shortcut cannot reproduce this because setting both $c_i$ to the floor gives
$\mBH_1(1-\Brat^{k})=\mBH_1(1-e^{-kL})>0$ for every $0<\Brat<1$. It
contains no image positions and has a fixed sign, so it cannot track a
sign that is set by image geometry.
\emph{[Attested: \solnlink{21_Perturbative_Estimators/problems\_21.wl}{problems\_21.wl}, checks 4a--4e.]}

% ----- Exercise: coherence limit -----
\subsection*{Exercise~\ref{ex:pe_coherence_limit}: Coherence Limit}

From $r_i=(\gp-1)[c_1-\Brat_i^{k}c_i]$ with $k>0$, $\Brat_i^{k}\to0$ as
$\Brat_i\to0$, so
\[
  \lim_{\Brat_i\to0}r_i=(\gp-1)\,c_1,
\]
independent of $c_i$ and of which plane was chosen.

In this limit the central-mass fraction seen by
plane $i$, $\mBH_i=\mBH_1\Brat_i^{k}$, tends to zero: plane $i$ is
effectively unperturbed and contributes nothing of its own. What
survives is plane~1's error. The distant plane acts as a clean
reference, and the ratio measures plane~1's contamination in full.
Adding such planes does not dilute the systematic; every ratio inherits
the same $c_1$.

As to how closely the limit can be approached, since $\Brat_i=d_1/d_i$ and every distance factor
$d_i=D_{{\rm d}s_i}/D_{s_i}<1$, we have $\Brat_i>d_1$. Even a source at
infinite redshift leaves $\Brat_i$ above $d_1$, so the limit is
approached but not reached; it is approached fastest for shallow slopes,
where $k$ is large. For equal $c_i=c>0$, $|r_i|$ increases
monotonically as $\Brat_i$ decreases.
\emph{[Attested: \solnlink{21_Perturbative_Estimators/problems\_21.wl}{problems\_21.wl}, checks 5a--5d.]}

% ----- Exercise: design -----
\subsection*{Exercise~\ref{ex:pe_design}: Design}

Neither move improves accuracy; both improve precision. Halving the
per-ratio error: for $C=\sigma^2I$ the bias
$(J^{\mathsf T}C^{-1}J)^{-1}J^{\mathsf T}C^{-1}\boldsymbol\delta$
does not depend on $\sigma$, while the marginal errors scale as
$\sigma$, so halving $\sigma$ leaves the bias unchanged and halves the
error, and $\mathcal{S}$ doubles. Adding a plane: in the
equal-sensitivity toy with a coherent shift,
$\mathcal{S}_N=\sqrt N\,\mathcal{S}_1$, so going from three to four
ratios raises $\mathcal{S}$ by $\sqrt{4/3}\approx1.15$.
Both make the result more precisely wrong. The quantity to
optimise is $\mathcal{S}$, not the statistical error, and the only lever
that lowers $\mathcal{S}$ is the modelling one: give the deflector's
inner profile enough freedom to represent a central mass, or resolve the
radial image well enough to constrain it.

One caveat is needed. A new plane adds an
independent row to the Jacobian and can break parameter degeneracies,
which is a genuine gain that the toy model ignores. The correct
statement is that extra planes buy constraining power but no
protection: unless the modelling systematic is addressed first,
the added power is spent sharpening a biased answer.
\emph{[Attested: \solnlink{21_Perturbative_Estimators/problems\_21.wl}{problems\_21.wl}, checks 6a--6d; the degeneracy-breaking caveat is conceptual.]}

% ----- Exercise: mass sheet -----
\subsection*{Exercise~\ref{ex:pe_mass_sheet}: A Different Omitted Component}

(a) A uniform sheet of convergence $\kappa$ deflects by
$\kappa\theta$, so the true source is
$\ysrc=\theta_j-\alpha_{\rm PL}(\theta_j)-\kappa\theta_j$. Repeating the
subtraction of \S\ref{sec:pe_residuals} gives
$r_j=\kappa a_j-\epsilon b_j-\delta\ysrc$ with $a_j=\theta_j$.

(b) For the canonical triple $a=(1,-1,\rho)$, so
$\bar a=\rho/3$, $\sum_ja_jb_j=(\gp-1)[2+\rho^{3-\gp}]$ and
$\operatorname{Cov}(a,b)=2(\gp-1)[1+\tfrac13\rho^{3-\gp}]$; the variance
is unchanged. Hence
\[
  c_{\rm sheet}=\frac{3+\rho^{3-\gp}}{(\gp-1)\bigl(3+\rho^{4-2\gp}\bigr)}.
\]

(c) Both $\rho^{3-\gp}$ and $\rho^{4-2\gp}$ vanish as $\rho\to0$,
so $c_{\rm sheet}\to1/(\gp-1)$; the tangential pair alone also gives
$1/(\gp-1)$ (there $a=b/(\gp-1)$ again). Solving the total lens equation
at first order gives $\thetaE^{\rm tot}=\thetaE[1+\kappa/(\gp-1)]$, the
same floor: the pair measures the total enclosed mass whatever its
form. At $\rho=0.05$, $\gp=1.6$: $c_{\rm sheet}=1.626$, slightly
below the floor $1.667$, while the point-mass $c=4.87$.

The reason is that leverage comes from where the omitted shape and the
model's direction differ most. A point mass's deflection $1/\theta$ is
largest at the centre, exactly where the power law's $b$ vanishes. A
sheet's deflection $\kappa\theta$ vanishes at the centre too, so
at the radial image both shapes are small and the radial image carries
almost no information about a sheet. Leverage is a property of the pair
(omitted component, datum), not of the datum alone.
\emph{[Attested: \solnlink{21_Perturbative_Estimators/problems\_21.wl}{problems\_21.wl}, checks 7a--7f.]}
